\documentclass[12pt]{article}
\usepackage[utf8]{inputenc}
\usepackage[T1]{fontenc}
\usepackage{hyperref}
\title{Yupik, Sirenik Language Resources for K-12 Curriculum Development}
\author{Compiled by Claude Sonnet 5.5 (claude-sonnet-5-5) for Andr\'as Kornai}
\date{}
\begin{document}
\maketitle

\section{Introduction}
Sirenik Yupik (ISO 639-3 \texttt{ysr}, Glottocode sire1246) is an Eskimo-Aleut language formerly spoken at Sireniki on the Chukchi Peninsula, Russia. It is regarded as extinct. Sources differ on when the last speaker died. One says January 1997, and another notes that censuses as late as 2010 reported up to 5 native speakers. Scholars also disagree about its classification. Some treat it as a third branch of Eskimo, and others as a Yupik language. Any curriculum would be a heritage and revitalization curriculum. The search found no school materials.

\section{Primary Educational Resources}
No primers, textbooks or school programmes for Sirenik were found.

\section{Linguistic and Reference Materials}
\begin{itemize}
\item Wikipedia, Sirenik language. It cites Menovshchikov's 1964 \emph{Jazyk sirenikskix eskimosov} (Moscow: Nauka), which describes the language in Cyrillic transcription: \url{https://en.wikipedia.org/wiki/Sirenik_language}
\item Max Planck EVA numeral-systems page for Sirenik Yupik, whose data source is Menovshchikov: \url{https://lingweb.eva.mpg.de/channumerals/Yupik-Sirenik.htm}
\item \'Etudes/Inuit/Studies article (PDF) that surfaced in the search; I did not verify its content or relevance: \url{https://www.erudit.org/fr/revues/etudinuit/2005-v29-n1-2-etudinuit1430/013937ar.pdf}
\end{itemize}

\section{Religious and Community Resources}
None found.

\section{Digital and Audio Resources}
\begin{itemize}
\item Endangered Languages Project page for ysr: \url{https://www.endangeredlanguages.com/lang/ysr?hl=en}
\end{itemize}
No audio recordings or corpus were found.

\section{Curriculum Development Context}
Revitalization would depend on the archival description, mainly Menovshchikov's work. This search found no dictionary, so a lexicon would have to be built from the archival material. Any work should involve descendant communities, who may be Siberian Yupik or Chukchi speakers. Related living languages such as Siberian Yupik could supply context, but they are different languages.
\end{document}